\documentclass[12pt]{article}
\usepackage{fullpage}
\usepackage{amsmath,amssymb,amsthm}

\newtheorem{lemma}{Lemma}
\newtheorem{proposition}{Proposition}
\newcommand{\R}{\mathbb R}
\newcommand{\Z}{\mathbb Z}
\newcommand{\ellg}{\operatorname{length}}

\begin{document}

\begin{center}
{\bf The infimum is \(\sqrt3\)}
\end{center}

Let
\[
        M_\beta=\Sigma/\langle G_\beta\rangle .
\]
This is the flat \(1\times \beta\) Mobius band.  If \(\beta\) is
realized, the given map descends, by condition (2), to an embedding
\(F:M_\beta\to \R^3\).  The invariant triangulation descends to a finite
clean triangulation of \(M_\beta\).

In a clean triangle, call the boundary edge the ridge and the opposite
vertex the apex.  A bend is a line segment in the triangle joining the
apex to a point of the ridge.

\begin{lemma}[all bends are lengthened]\label{lem:lengthen}
Let \(L\) be the affine map on one clean triangle.  If \(P\) is any point
of the ridge and \(C\) is the apex, then
\[
        |L(P)-L(C)|>|P-C|.
\]
Also every non-degenerate subarc of the boundary is shortened by \(F\).
\end{lemma}

\begin{proof}
Let the ridge endpoints be \(A,B\), and put
\(P_s=(1-s)A+sB\).  The function
\[
 q(s)= |L(P_s)-L(C)|^2-|P_s-C|^2
\]
is quadratic and
\[
 q''(s)=2\big(|L(A)-L(B)|^2-|A-B|^2\big)<0.
\]
Thus \(q\) is concave.  The two non-ridge sides are lengthened, so
\(q(0)>0\) and \(q(1)>0\); concavity gives \(q(s)>0\) for
\(0\le s\le1\).  On a boundary edge the affine map has constant Lipschitz
constant \(<1\), so every boundary subarc is shortened.
\end{proof}

\begin{lemma}[a polygonal \(T\)-pair]\label{lem:T}
For the embedded polygonal Mobius band \(F(M_\beta)\) there are two bends
whose image segments lie on perpendicular intersecting lines.
\end{lemma}

\begin{proof}
The bends form a circle.  Indeed, in each triangle the segments from the
apex to the ridge foliate the triangle.  At a common non-boundary edge
the two adjacent one-parameter families have the same limiting bend, so
the intervals of bends glue together.  Equivalently, every bend meets the
center circle \(\{y=1/2\}/\langle G_\beta\rangle\) exactly once, and every
point of this center circle lies on exactly one bend.

Parametrize this bend circle by \(u\in\R/\Z\), and lift to \(u\in\R\).
Choose continuously an orientation of the lifted bends from the lower
boundary line to the upper one.  If \(e(u)\) is the corresponding unit
direction vector of the image segment and \(m(u)\) its midpoint, then
\[
        e(u+1)=-e(u),\qquad  m(u+1)=m(u),
\]
because going once around the Mobius band applies \(G_\beta\), which
interchanges the two boundary components.

For \(0<r<1\) define
\[
 \Phi(u,r)=\Big(e(u)\!\cdot e(u+r),\,
 (m(u)-m(u+r))\!\cdot (e(u)\times e(u+r))\Big).
\]
This is well-defined on an open cylinder.  It extends continuously to the
two-point compactification of that cylinder, a \(2\)-sphere, with limiting
values \((1,0)\) as \(r\to0\) and \((-1,0)\) as \(r\to1\).  The involution
\((u,r)\mapsto(u+r,1-r)\), which swaps the two bends, is conjugate to the
antipodal map on this sphere, and the above identities give
\(\Phi\circ\iota=-\Phi\).  By the Borsuk--Ulam theorem, \(\Phi=0\) at
some interior point.

At such a point the two direction vectors are perpendicular.  The second
coordinate says that the vector joining the two midpoints lies in the
plane spanned by these directions; hence the two supporting lines are
coplanar.  Since they are not parallel, they intersect.
\end{proof}

\begin{proposition}\label{prop:lower}
Every realized \(\beta\) satisfies \(\beta>\sqrt3\).
\end{proposition}

\begin{proof}
Choose the two bends from Lemma \ref{lem:T}.  Their image segments have
disjoint relative interiors, because distinct bends have disjoint
relative interiors in \(M_\beta\) and \(F\) is an embedding.  Let \(O\) be
the intersection of the two supporting lines.  At least one of the two
image segments does not contain \(O\) in its relative interior; call that
bend \(B\), and call the other one \(T\).

Lift \(T\) to a segment from \((0,0)\) to \((t,1)\), after reflecting the
\(x\)-axis if necessary, with \(0\le t<\beta\).  (If the horizontal
displacement had absolute value at least \(\beta\), the lift would meet a
\(G_\beta\)-translate of itself.)  Cutting \(M_\beta\) along
\(T\) gives two boundary arcs \(H,D\) of lengths
\[
        \ellg(H)=\beta-t,\qquad \ellg(D)=\beta+t.        \tag{1}
\]
Let \(T^*=F(T)\).  Since \(H^*=F(H)\) joins the endpoints of \(T^*\), and
boundary arcs are strictly shortened while bends are lengthened,
\[
        \beta-t=\ellg(H)>\ellg(H^*)\ge |T^*|>\sqrt{1+t^2}.       \tag{2}
\]

The segment \(B^*=F(B)\) lies on a line perpendicular to the line of
\(T^*\).  Since \(O\) is not in the relative interior of \(B^*\), one
endpoint \(v^*\) of \(B^*\) has distance from the line of \(T^*\) at least
\(|B^*|\).  By Lemma \ref{lem:lengthen}, \(|B^*|>1\).  Let \(A\) be the
one of the two boundary arcs \(H,D\) which contains the corresponding
boundary point \(v\).  The image arc \(A^*\) goes from one endpoint of
\(T^*\) to the other through \(v^*\).  If \(V\) denotes the sum of the two
straight distances from \(v^*\) to the endpoints of \(T^*\), then
\(\ellg(A^*)\ge V\).  A triangle with base \(|T^*|>\sqrt{1+t^2}\) and
height \(>1\) has the sum of its two non-base sides strictly bigger than
\[
        \sqrt{(1+t^2)+4}=\sqrt{5+t^2}.
\]
Therefore
\[
        \ellg(A)>\ellg(A^*)\ge V>\sqrt{5+t^2}.                  \tag{3}
\]

If \(A=H\), then \(D\ge H\) by (1), and (3) gives
\[
        2\beta=H+D>2\sqrt{5+t^2}>2\sqrt3.
\]
So assume \(A=D\).  Combining (1), (2), and (3), we get
\[
 2\beta=H+D>\sqrt{1+t^2}+\sqrt{5+t^2},                         \tag{4}
\]
and also
\[
 2\beta=2D-2t>2\sqrt{5+t^2}-2t.                                \tag{5}
\]
Set
\[
 a(t)=\sqrt{1+t^2}+\sqrt{5+t^2},\qquad
 b(t)=2\sqrt{5+t^2}-2t.
\]
The function \(a\) is increasing on \([0,\infty)\), while \(b\) is
decreasing there, and at \(t_0=1/\sqrt3\) one has
\[
        a(t_0)=b(t_0)=2\sqrt3 .
\]
Thus \(\max(a(t),b(t))\ge2\sqrt3\) for every \(t\ge0\).  Equations
(4)--(5) imply \(2\beta>2\sqrt3\).
\end{proof}

It remains to show that numbers arbitrarily close to \(\sqrt3\) from
above are realized.

\begin{lemma}[realizations above \(\sqrt3\)]\label{lem:upper}
Every \(\beta>\sqrt3\) is realized.
\end{lemma}

\begin{proof}
We use the classical triangular construction: for every
\(\lambda>\sqrt3\) there is an embedded polygonal paper Mobius band
\(P:M_\lambda\to\R^3\), affine and isometric on the triangles of a clean
triangulation.  This is the usual upper-bound construction obtained from
the triangular Mobius band by making the strip slightly longer and
separating (or rounding, then polyhedrally approximating) the three
folds.

Choose \(\lambda\) with \(\sqrt3<\lambda<\beta\), and put
\(s=\lambda/\beta<1\).  The map
\[
        h:M_\beta\longrightarrow M_\lambda,\qquad h(x,y)=(sx,y),
\]
is a homeomorphism because \(h\circ G_\beta=G_\lambda\circ h\).
Pull the clean triangulation of \(M_\lambda\) back by \(h\).  There are
only finitely many edge orbits.  For each non-boundary edge \(e\) of the
pulled-back triangulation, let \(d_e\) be its horizontal displacement in
the \(M_\beta\) coordinates.  Since the set of such \(e\) is finite, we
may choose \(r\) with
\[
        1<r<1/s
\]
so close to \(1/s\) that, for every non-boundary edge \(e\),
\[
        r\sqrt{1+s^2d_e^2}>\sqrt{1+d_e^2}.                    \tag{6}
\]
Now define \(F:M_\beta\to\R^3\) by \(F=r\,P\circ h\), and lift it to a
map \(f:\Sigma\to\R^3\).  The map is affine on every pulled-back triangle
and is invariant under \(G_\beta\).  Since \(h\) is a homeomorphism of
quotients and \(P\) is an embedding, the lifted map satisfies
\(f(p)=f(q)\) if and only if \(p=G_\beta^k(q)\) for some \(k\in\Z\).

On a boundary edge of length \(L\) in \(M_\beta\), the image length is
\(rsL<L\).  On a non-boundary edge with displacement \(d_e\), the original
length is \(\sqrt{1+d_e^2}\), while the image length is
\(r\sqrt{1+s^2d_e^2}\), which is larger by (6).  Thus the restriction to
each triangle is a squeeze map.
\end{proof}

Proposition \ref{prop:lower} shows that no realized number is
\(\le\sqrt3\), and Lemma \ref{lem:upper} shows that every
\(\beta>\sqrt3\) is realized.  Hence the infimum of the set of realized
numbers is \(\sqrt3\).

\end{document}
